\section{Models, Assumptions, Equations, and Interpretation}

An equation is never the whole theory by itself. A field equation becomes a
physical model only after one specifies the field, the domain, the assumptions,
the interpretation of the variables, and the class of allowed solutions.

For instance, the wave equation
\[
    \frac{\partial^2 \phi}{\partial t^2}
    -
    c^2 \Delta \phi
    =
    0
\]
may describe small vibrations of a string, sound waves in an idealized medium,
or a simple relativistic scalar field after a change of notation. The same
formal equation can therefore represent different physics. The interpretation
comes from what \(\phi\) means, what \(c\) means, what the domain is, and what
approximations led to the equation.

Throughout these notes, a model will be read through four layers:
\[
\text{assumptions}
\longrightarrow
\text{mathematical variables}
\longrightarrow
\text{field equations}
\longrightarrow
\text{physical interpretation}.
\]
This order matters. If the assumptions are changed, the same equation may no
longer be appropriate. If the interpretation is changed, the same symbols may
refer to different observables.

\begin{summarybox}
The goal is not only to manipulate field equations. The goal is to understand
what each equation says, what it ignores, which quantities it conserves, and
which physical questions it can answer.
\end{summarybox}
